% !TEX program = xelatex
\documentclass{英语试卷}
\begin{document}
\examtitle{浙江强基联盟 2025 年 5 月高三联考}{英语}
\par 浙江强基联盟 2025 年 5 月高三联考\par
\par 英语学科 试题卷\par
\par 浙江强基联盟研究院 命制\par
\begin{notice}
\par 1. 答卷前，考生务必将自己的姓名、准考证号填写在答题卡上。\par
\par 2. 回答选择题时，选出每小题答案后，用铅笔把答题卡上对应题目的答案标号涂黑。如需改动，用橡皮擦干净后，再选涂其他答案标号。回答非选择题时，将答案写在答题卡上，写在本试卷上无效。\par
\par 3. 考试结束后，将本试卷和答题卡一并交回。\par
\end{notice}
\parttitle{第一部分 听力（共两节，满分 30 分）}
\par 做题时，先将答案标在试卷上，录音内容结束后，你将有两分钟的时间将试卷上的答案转涂到答题卡上。\par
\sectitle{第一节（共5小题：每小题1.5分，满分7.5分）}
\par 听下面 5 段对话。每段对话后有一个小题，从题中所给的 A、B、C 三个选项中选出最佳选项。听完每段对话后，你都有 10 秒钟的时间来回答有关小题和阅读下一小题。每段对话仅读一遍。\par
\question{1. What kind of food will the man have with his friends next month?}
\begin{choices}(3)
  \item Mexican food.
  \item Italian food.
  \item French food.
\end{choices}
\question{2. What will the weather probably be like this weekend?}
\begin{choices}(3)
  \item Sunny.
  \item Rainy.
  \item Hot.
\end{choices}
\question{3. What will the woman do?}
\begin{choices}(3)
  \item Have a rest at home.
  \item Broadcast live on social media.
  \item Go to the poetry reading event.
\end{choices}
\question{4. What is the probable relationship between the speakers?}
\begin{choices}(3)
  \item Friends.
  \item Boss and employee.
  \item Interviewer and candidate.
\end{choices}
\question{5. Who will take care of the dogs?}
\begin{choices}(3)
  \item The woman's niece.
  \item The woman's sister.
  \item The woman's parents.
\end{choices}
\sectitle{第二节（共15小题：每小题1.5分，满分22.5分）}
\par 听下面 5 段对话或独白。每段对话或独白后有几个小题，从题中所给的 A、B、C 三个选项中选出最佳选项。听每段对话或独白前，你将有时间阅读各个小题，每小题 5 秒钟；听完后，各小题将给出 5 秒钟的作答时间。每段对话或独白读两遍。\par
\par 听第 6 段材料，回答第 6、7 题。\par
\question{6. Which character would the woman like to be?}
\begin{choices}(3)
  \item Juliet.
  \item Rosa.
  \item Nurse.
\end{choices}
\question{7. What does the woman care most about the play?}
\begin{choices}(3)
  \item Performing well.
  \item Playing the lead role.
  \item Learning about the story.
\end{choices}
\par 听第 7 段材料，回答第 8 至 10 题，\par
\question{8. What happened to the woman?}
\begin{choices}(3)
  \item She lost her bus pass.
  \item She took the wrong bus.
  \item She missed the school bus.
\end{choices}
\question{9. What day is it today?}
\begin{choices}(3)
  \item Wednesday.
  \item Thursday.
  \item Friday.
\end{choices}
\question{10. What will the woman probably do after getting to school?}
\begin{choices}(3)
  \item Look for Mr. Smith
  \item Borrow some money.
  \item Go to the front office.
\end{choices}
\par 听第 8 段材料，回答第 11 至 13 题：\par
\question{11. What does the woman say about the book for her report?}
\begin{choices}(3)
  \item It's lost.
  \item It's damaged.
  \item It's sold out.
\end{choices}
\question{12. Why is the man's uncle mentioned in the conversation?}
\begin{choices}(3)
  \item He is a famous writer.
  \item He used to work in the library.
  \item He has many books on architecture.
\end{choices}
\question{13. What will the man probably do next?}
\begin{choices}(3)
  \item Go to the library.
  \item Attend a class.
  \item Make a call.
\end{choices}
\par 听第 9 段材料，回答第 14 至 17 题。\par
\question{14. What is similar between math problems and making eggs?}
\begin{choices}(3)
  \item They're both fun and easy to deal with.
  \item They're both best to do in the morning.
  \item There are many different ways to do them.
\end{choices}
\question{15. What kind of eggs does the boy like best?}
\begin{choices}(3)
  \item Scrambled ones.
  \item Boiled ones.
  \item Fried ones.
\end{choices}
\question{16. What will the boy do next?}
\begin{choices}(3)
  \item Continue with his homework.
  \item Get help from Mr. Thompson.
  \item Eat dinner upstairs in his room.
\end{choices}
\question{17. Where are the speakers talking?}
\begin{choices}(3)
  \item In a classroom.
  \item At home.
  \item In an office.
\end{choices}
\par 听第 10 段材料，回答第 18 至 20 题。\par
\question{18. Who controlled Saudi Arabia in the 19th century?}
\begin{choices}(3)
  \item The British.
  \item The French.
  \item The Ottoman Empire.
\end{choices}
\question{19. When was the independent kingdom of Saudi Arabia founded?}
\begin{choices}(3)
  \item In 1916.
  \item In 1932.
  \item In 1938.
\end{choices}
\question{20. What happened right after the Saudis discovered oil?}
\begin{choices}(3)
  \item They made investments in finance.
  \item They signed agreements with other nations.
  \item They formed a brand new industry on their own.
\end{choices}
\parttitle{第二部分 阅读理解（共两节，满分 50 分）}
\sectitle{第一节（共15小题：每小题2.5分，满分37.5分）}
\par 阅读下列短文，从每题所给的 A、B、C 和 D 四个选项中，选出最佳答案，\par
\passagelabel{A}
\par The Best Online Flower Delivery Services\par
\par It's easy to tie yourself in knots trying to find the perfect Mother's Day gift. Flowers are a very convenient present. Many online flower delivery services offer same-or next-day delivery.\par
\par Marks \& Spencer\par
\par Marks \& Spencer has recently improved its flower business, offering a five-day freshness guarantee on its bouquets. M\&S sources as many British flowers as it can, depending on what's in season. It also uses sustainable packaging.\par
\par Fastest delivery: next day\par
\par Delivery cost: free Scilly Flowers\par
\par All Scilly Flowers stems are grown naturally outdoors without artificial light or heat. As such, these particular flowers are available from October to Easter. The whole farm is making strides in sustainability with low-irrigation initiatives, a solar power cold store and postal plane and boat shipping.\par
\par Fastest delivery: 1st class Royal Mail Delivery cost: free Smith \& Munson\par
\par Smith \& Munson grows all its flowers in Lincolnshire-it specialises in tulips, lilies, peonies and ranunculus. These are all seasonal, so you can't get them out of season. They are all buds when they arrive, ensuring they last as long as possible.\par
\par Fastest delivery: 24 hours (Tuesday — Saturday) Delivery cost: free Interflora\par
\par Interflora is effectively a flower matchmaker. You choose from a huge menu of flowers and then it matches you with a florist(花店) local to your delivery address. The florist then creates and delivers the arrangement.\par
\par Fastest delivery: same day(order before 3pm)\par
\par Delivery cost: £10 for same-day delivery, £7.50 for named-day delivery\par
\question{21. What is special about Marks \& Spencer?}
\begin{choices}(2)
  \item It grows seasonal flowers.
  \item It guarantees a week's freshness.
  \item It offers same-day delivery.
  \item It uses nature-friendly wrapping
\end{choices}
\question{22. Which flower delivery service charges for delivery?}
\begin{choices}(2)
  \item Marks \& Spencer.
  \item Interflora.
  \item Smith \& Munson.
  \item Scilly Flowers.
\end{choices}
\question{23. Which section of the newspaper is the text probably taken from?}
\begin{choices}(2)
  \item Art.
  \item Travel.
  \item Lifestyle.
  \item Health.
\end{choices}
\passagelabel{B}
\par As I wandered out of my New York apartment, behind me someone called out, "Where can I find peaches?" I turned around to see a woman standing outside the greengrocer's. The absurdity of the moment struck me—why would someone want peaches in the middle of winter? It is just as absurd as sitting by the pool on a blistering summer day and reaching for a warm, woolly jumper.\par
\par I was already aware of the issues facing the food system: industrial farming ruining soils, unclear supply chains leaving citizens powerless in making the right buying decisions, and the dominance of ultra-processed foods with zero nutritional value in supermarkets. But this encounter made me see how normalized it was to expect food out of season.\par
\par After graduating, I worked at Wall Street in New York, managing portfolios(投资组合). Despite the breathtaking view and fast-paced stock market, I was troubled by the greed underlying my work. Cooking was my escape. It had always been a passion, whether in Buenos Aires, at university, or in New York.\par
\par As I neared my 30s, I knew I needed a change and food became my north star. Through my work visiting markets and farmers around Europe, I saw first-hand why seasonality is critical in our ability to regain control of, and fix, our food system. All around me I could see how quickly flavour was disappearing from our plates. I wanted to reintroduce phenomenal flavour into people's lives, and to re-educate and reconnect consumers with nature. In 2004, I took over\par
\par Natoora, a French-founded company with a small base in London, determined to reintroduce great flavor to people's lives and reconnect them with nature.\par
\par My two-decade-long devotion to flavor, initially a pursuit of personal pleasure, has given my work a deep sense of purpose. Flavor can emotionally connect us to food and bring joy back into our relationship with it. Changing careers opened a new world for me, filled with purpose and love for my work, and I hope it inspires others to fight for a better food-centric world.\par
\question{24. Why does the author mention the encounter with a woman?}
\begin{choices}(2)
  \item To illustrate an absurd shopping choice.
  \item To highlight a problem in the food system.
  \item To compare between peaches and jumpers.
  \item To show people's separation from farming.
\end{choices}
\question{25. What does the author think of cooking?}
\begin{choices}(2)
  \item Healing.
  \item Troublesome.
  \item Demanding.
  \item Instructive.
\end{choices}
\question{26. What can we infer from paragraph 4?}
\begin{choices}(2)
  \item People will improve flavors on their tables soon.
  \item Consumers are not willing to get close to nature.
  \item Our food system is in need of our great attention.
  \item All the markets in Europe sell food out of season.
\end{choices}
\question{27. Which of the following can be a suitable title for the text?}
\begin{choices}(2)
  \item Add a Unique flavor to Your career
  \item The Great Power of Seasonal Eating
  \item Approaches to Fixing the Food System
  \item From Wall Street to the World of Flavor
\end{choices}
\passagelabel{C}
\par When babies are exposed to a variety of faces, they become more adept in recognizing and appreciating diversity, which can help reduce biases and prejudices as they grow. This is according to neuroscientist Saskia Koch of the Donders Institute and Radboud University.\par
\par Koch conducted her research using data from the Nijmegen longitudinal study. This study follows young people from infancy(婴儿期). First, the researchers looked at how frequently the teenagers had attended day care as babies. Then, the same youngsters were assigned the duty of playing a game when they reached the age of 17. In the game, they had to work with another person they couldn't see. Both players had the same goal: find an object on a gameboard. Only the 17-year-old knew where the item was. The 17-year-olds had to communicate this location to each other without using words. They had to design communicative patterns that would assist the other person in locating the object. This could be accomplished, for example, by performing specific movements.\par
\par The 17-year-olds were informed they were playing with a 5-year-old in one game and an adult in another, despite the fact that it was the same individual each time. Initially, the 17-year-olds adjusted their communicative patterns when they thought they were interacting with the child by placing more emphasis on their movements. How quickly they abandoned this pattern depended on how much time they spent in day care. 17-year-olds who had come into contact with many different people as infants were more likely to get over a stereotype(刻板印象). They found out in the interaction that the "child" was also very good and then based their play on that.\par
\par According to the researcher this means that babies who have had more diverse social interactions can more easily let go of a stereotype later in life. They are more sensitive to social interaction and can more quickly let go of the image they have in their minds of their peers. Koch says, "Gaining many diverse social experiences early in life allows us to adapt better to the needs of our partners in social interaction."\par
\question{28. What does the underlined word "adept" in paragraph 1 mean?}
\begin{choices}(2)
  \item Forgetful.
  \item Skillful.
  \item Interested.
  \item Disappointed.
\end{choices}
\question{29. What was the purpose of the game in the study?}
\begin{choices}(2)
  \item To confirm the data from the Nijmegen study.
  \item To test the impacts of early social interactions.
  \item To enhance the youngsters' cooperative ability.
  \item To assess the effects of communicative patterns.
\end{choices}
\question{30. How did the 17-year-olds first respond to the "child" in the game?}
\begin{choices}(2)
  \item They changed their fixed idea of a child.
  \item They adjusted the time spent on the game.
  \item They performed more specific movements.
  \item They employed verbal interactive patterns.
\end{choices}
\question{31. What helps babies with future social contacts according to Koch?}
\begin{choices}(2)
  \item Interacting with various people.
  \item Engaging with diverse peers.
  \item Adapting to the needs of partners.
  \item Letting go of the old images.
\end{choices}
\passagelabel{D}
\par Planetary geoarchaeology(地质考古学), the study of how cultural and natural processes on Earth's moon, on Mars, and throughout the solar system may be changing, preserving, or erasing the material record of space exploration, has been proposed by two researchers from the Kansas Geological Survey at the University of Kansas and their colleagues as a new scientific subfield.\par
\par "Until recently, we might consider the material left behind during the space race of the mid-20th century as relatively safe," said Justin Holcomb, researcher at the Kansas Geological Survey. "However, the material record that currently exists on the moon is rapidly becoming at risk of being destroyed if proper attention isn't paid during the new space era."\par
\par Applying geoarchaeological tools and methods to the movement of people into space and the solar system is a natural extension of the study of human migration on Earth, the focus of the ODYSSEY Archaeological Research Program directed by Holcomb's co-author, Rolfe Mandel.\par
\par How planetary geoarchaeologists will determine whether an item is worth preserving is an open question. "We feel that all material currently existing on extraterrestrial(地球外的)surfaces is space heritage and worthy of protection," Holcomb said. "However, some sites, such as the very first footprints on the moon represent the material footprint of a long history of migration." Beyond those "firsts", making decisions on a case-by-case basis will be necessary to sort through the tens of thousands of pieces of debris(残骸) now in orbit or dropped across the surfaces of the moon and Mars, which many people refer to as "trash" but Holcomb and his colleagues see as heritage.\par
\par Holcomb thinks geoarchaeologists ought to be a part of upcoming NASA missions to ensure the security and preservation of space history. As for being part of a space mission himself? "I'll leave that to other geoarchaeologists," Holcomb said. "There's plenty to do down here, but I do hope to see an archaeologist in space before it's all over."\par
\question{32. Why has planetary geoarchaeology been proposed?}
\begin{choices}(2)
  \item To protect space heritage.
  \item To study the solar system.
  \item To resist space exploration.
  \item To found a scientific field.
\end{choices}
\question{33. What do we know about the material record of space?}
\begin{choices}(2)
  \item It threatens the movement of people into space.
  \item It is the main focus of the ODYSSEY Program.
  \item It shows the whole process of human migration.
  \item It is not as safe as expected in the new space era.
\end{choices}
\question{34. What does paragraph 4 mainly talk about?}
\begin{choices}(2)
  \item The methods of classifying space trash.
  \item The first footprints on the moon and Mars.
  \item The question of planetary geoarchaeologists.
  \item The discussion on identifying space heritage.
\end{choices}
\question{35. What does Holcomb suggest other geoarchaeologists do?}
\begin{choices}(2)
  \item Learn from scientists like him.
  \item Lead upcoming NASA missions.
  \item Participate in space exploration.
  \item Ensure the security of astronauts.
\end{choices}
\sectitle{第二节（共5小题；每小题2.5分，满分12.5分）}
\par 阅读下面短文，从短文后的选项中选出能填入空白处的最佳选项。选项中有两项为多余选项，\par
\par How to Clean a House\par
\par Cleaning your house seems like a fairly self-explanatory task, but when you sit down to think about it, you might realize that you have a lot of questions. \blanknum{36} We'll work it down into a process that seems simple enough to tackle.\par
\par Decide how clean you want your house to be and how much time you have. \blanknum{37} So, be honest with yourself about what you can do, how much time you have, and how motivated you're feeling.\par
\par \blanknum{38} You don't want to vacuum and then get crumbs all over the floor, or dust the bottom of something and knock more dust down when you work your way up onto the previously uncleaned areas. If you don't have much time, start with the mess first and work your way to more serious things from there.\par
\par Make a checklist and a plan of attack. Know which rooms you plan to start and finish your cleaning in. This will help speed up the process and ensure you are not repeating your efforts, especially if more than one person is involved in the cleaning. You can follow the steps in the cleaning part in order. \blanknum{39}\par
\par Attempt to delegate tasks. If you live with other people, cleaning the house should not be exclusively your responsibility! You may have to take on leadership by setting up a rotational(轮流的) cleaning schedule, but it's far better than having to do this backbreaking work alone. \blanknum{40} For example, grade school children can pick up their bedroom floors, teenagers can clean the garage or bathroom, etc.\par
\par Once you get started, you won't want to stop until the whole house is spotless and organized.\par
\begin{choices}(1)
  \item For example, how do you clean a toilet?
  \item Work from the top down, if at all possible.
  \item Once you get started, make sure it is equal.
  \item Make sure the tasks assigned are age-appropriate.
  \item Should this really be something you do in everyday life?
  \item This way, you'll have a checklist already written out for you.
  \item This wil help you decide how to set up your cleaning schedule.
\end{choices}
\parttitle{第三部分 语言知识运用（共两节，满分 30 分）}
\sectitle{第一节（共15小题；每小题1分，满分15分）}
\par 阅读下面短文，从短文后各题所给的 A、B、C 和 D 四个选项中，选出可以填入空白处的最佳选项。\par
\par It started like any normal day. My partner, Stewart, helped me out of bed. I \blanknum{41} on my crutches(拐杖) as I made my way to the car, slowly and dangerously. All of a sudden, I \blanknum{42} my footing and fell.\par
\par Not for the first time, Stewart was caught between concern and \blanknum{43} Wasn't it time, he asked, for me to consider using a wheelchair? Once again, I \blanknum{44} the suggestion.\par
\par I was diagnosed with rheumatoid arthritis(类风湿性关节炎) in my 20s. Over the following years, movement and mobility became more and more painful and challenging, but I was \blanknum{45} to stay on my feet — to still be a "walking person". But, slowly, walking unaided became impossible. The morning after my \blanknum{46}, the truth of what Stewart had been saying \blanknum{47} ne: trying to stay on my feet was disabling me further. My perspective \blanknum{48} and I admitted what felt, at the time, like defeat. At 46, I started using the wheelchair. Suddenly, having been \blanknum{49} here I was, a "disabled person", sitting at the height of most people’s waists.\par
\par However, being more \blanknum{50} disabled has had other silver linings, too. People are kind. And, I often feel as if I am seeing the very best of human \blanknum{51}.\par
\par Now, when I look at the wheelchair, I see the imprint of myself on it, how its shape has \blanknum{52} me, and mine to it, and I feel greater \blanknum{53}: it is a part of me now.\par
\par Of course, I long to walk, but I am learning to live with my \blanknum{54}. Using a wheelchair can be challenging, but it has also proved hugely \blanknum{55} together, it feels as if we are starting a whole new adventure.\par
\cz{41}{swung}{lived}{leaned}{hung}
\cz{42}{missed}{treated}{gained}{enjoyed}
\cz{43}{happiness}{frustration}{guilt}{relief}
\cz{44}{received}{considered}{appreciated}{rejected}
\cz{45}{unwilling}{determined}{frightened}{hesitant}
\cz{46}{fall}{success}{journey}{walk}
\cz{47}{excited}{escaped}{hit}{beat}
\cz{48}{remained}{shifted}{disappeared}{narrowed}
\cz{49}{bright}{reliable}{mean}{tall}
\cz{50}{mentally}{completely}{mildly}{visibly}
\cz{51}{history}{nature}{achievement}{activity}
\cz{52}{referred to}{belonged to}{adapted to}{turned to}
\cz{53}{acceptance}{burden}{distance}{objective}
\cz{54}{hobby}{dream}{disability}{talent}
\cz{55}{dominating}{amusing}{disturbing}{liberating}
\sectitle{第二节（共10小题；每小题1.5分，满分15分）}
\par 阅读下面短文，在空白处填入 1 个适当的单词或括号内单词的正确形式。\par
\par Known as the "male embroiderer", Ji Qiyu is the seventh-generation inheritor of the Zhuang brocade(锦) technique. He has introduced innovations to the art form and has turned the village into \blanknum{} \blanknum{56}\blanknum{} fascinating tourist destination.\par
\par Now, about five tourist groups visit the village every month. These \blanknum{} visitor) spend their time sightseeing, relaxing, and attending cultural workshops, immersing themselves in the charm of the rural \blanknum{} (mountain) village and the unique Zhuang ethnic culture.\par
\par Ji \blanknum{59} (grow) up in a family of Zhuang brocade craftsmen. After graduation, Ji returned to his hometown \blanknum{60} founded a Zhuang brocade cooperative. His technological advances lowered the barrier for others to learn the craft, \blanknum{61} (offer) local women the opportunity to work from home and earn an income.\par
\par Ji's designs, \blanknum{62} combine ethnic costumes and brocade elements, have gained market recognition for their uniqueness and craftsmanship. His innovative approach has not only attracted attention but also \blanknum{63}(significant) boosted local job opportunities and sales of agricultural products.\par
\par Thanks to Ji's efforts, Zhuang brocade, known for its vibrant colors and intricate patterns, \blanknum{64} (preserve) for decades. With over a thousand years of history, the art form continues to be passed down from generation \blanknum{65} generation.\par
\parttitle{第四部分 写作（共两节，满分 40 分）}
\sectitle{第一节（满分 15 分）}
\par 假定你是李华，你校交换生 Alex 计划本周六独自去徒步探险，向你咨询相关信息，请你写一封回信，内容包括：\par
\par (1) 推荐一处目的地；\par
\par (2) 提出安全建议。\par
\par 注意:\par
\par (1) 写作词数应为 80 左右；(2) 可适当增加细节，以使行文连贯，\par
\par\vspace{0.4em}
\begin{center}\begin{tabular}{{|l|}}
\hline
Dear Alex, Yours, Li Hua \\
\hline
\end{tabular}\end{center}
\vspace{0.4em}
\sectitle{第二节（满分25分）}
\par 阅读下面材料，根据其内容和所给段落开头语续写两段，使之构成一篇完整的短文，\par
\par There's an image of big cities as soulless places, where everyone appears to be looking out only for themselves. However, my own experiences tell a different story. During some of the most challenging times of my life, when I was struggling to find proper housing, I worked as an artist on the street. For years, I played my flute(长笛) there, and in doing so, I discovered the natural kindness of people. They showed compassion without expecting anything in return, driven purely by human goodwill\par
\par This profound realization took hold when I was 21, living in Sydney. I worked at a local pub, feeling lost and directionless. Amid the chaos of my life, playing the flute was my only constant source of comfort. Looking back, my connection with the flute began when I was 16.Earning a mere \$2.35 an hour at the supermarket, I saved for two long years to buy it. Every day after work, playing the flute became my way of processing the day's emotions.\par
\par One day, while working at the pub, someone asked me what I was doing that night. "I'm going to be playing the flute at —" I started to say, but then it struck me: I had forgotten to pick up my flute when I got off the train. Panic set in immediately. In a fit of frustration, I threw the plastic chair I was holding across the room.\par
\par I rushed to the train station, hoping someone had handed in my flute. No luck. Still, I left my number, clinging to the slim hope that it would turn up. Convinced that it was gone forever — after all, I hadn't even written my name on the case—I spent the next week drowning in self-blame, drinking heavily and feeling sorry for myself.\par
\par Then, out of the blue, I received a call from CityRail. They had my flute. I hurried to the lost property office at Central Station. And there it was, among all the other lost items.\par
\par 注意:\par
\par (1) 续写词数应为 150 左右;\par
\par (2) 请按如下格式在答题卡的相应位置作答。\par
\writingpara{Paragraph 1: As I went to grab it, the guards there asked for my ID. \blanknum{}}
\writinglines{8}
\writingpara{Paragraph 2: Desperate and anxious, I said, "What if I play you a tune?"}
\writinglines{8}
\end{document}
